\scanpage{87}
% archive foliation: 89
\uline{NB} On sait que la donnée d'un octaèdre \inserted{comb.} non orienté
revient à celle d'un système $(I,\omega)$, \deleted{\unsure{dans}} $\operatorname{Ens}_4\times\operatorname{Ens}_2$,
ens des \unsure{octo-quartes} et des orientations. Ici $\omega=\{i,-i\}$, \deleted{\unsure{et on}}
\deleted{par la donnée de $I$, que}\inserted{Mais ici on n'a pas $I$ qui est \unsure{abstrait}, mais}
$S$, qui est un ens de cardinal 6 sur lequel $V$ opère.
La condition à mettre \deleted{est} pour que ceci définisse une
\deleted{des} structure octaédrale (dont $S$ sera l'ens des sommets)
est que $\forall\sigma\in V^*$, \deleted{$\unsure{V}$} $\operatorname{card}S^\sigma=2$, et $S=\bigcup_{\sigma\in V^*}S^\sigma$
(réunion forcément disjointe) (cela entraîne que
$\sigma$ opère sur $(S\setminus S^\sigma)$ par transposition de \deleted{\unsure{les}}\inserted{\unsure{chacune}} des deux
composantes $S^{\sigma'},S^{\sigma''}$, ici $V^*=\{\sigma,\sigma',\sigma''\}$). Donc la
donnée de \deleted{$\unsure{S,V\subset}$} $(V,S,u:V\hookrightarrow\unsure{\mathfrak S}_S)$, ou $(S,V\subset\unsure{\mathfrak S}_S)$,
équivaut à celle de $S$, avec une involution sans pts fixes
$a$ (ayant les $S^\sigma$ comme orbites).

Revenant à la donnée d'un $V$ qui opère \inserted{fid.\textsuperscript{t}} sur $X$,
considérons une orbite \inserted{$I$} de $V$ sur $X^*=X\setminus S(V)$.
Donc $I$ est un torseur sous $V$. La donnée de
$(X,I,I\hookrightarrow X)$ équivaut à celle de $(X,I,V(I)\hookrightarrow G_X,$
$I$ orbite de $V(I)$ dans $X)$, et on se propose de
comparer le $(V=V(I))$-torseur $I$, avec le
$V$-torseur $I(V)$ \inserted{$=I(V,X)$} défini via l'octaèdre $\operatorname{Oc}(V)=\operatorname{Oc}(V,X)$.

En termes de \inserted{$I\deleted{\unsure{V}},\{\pm i\}$}, on peut définir un octaèdre \deleted{\uline{orienté}}
canonique \deleted{\unsure{sans base marquée, défini par la}}
condition que \deleted{\unsure{le groupe de ses} \illegible{1} \unsure{soit} $V$} \inserted{l'ens de ses \unsure{2|8} \unsure{faces} ``\unsure{ant.}'' $I_\pm$, l'ens de ses orientations $\{\pm i\}$}
\deleted{\unsure{bipartitions}}.
Cet \deleted{oct.} octaèdre, \deleted{une fois choisi un $i=\sqrt{\unsure{1|-1}}$,} peut
se réaliser canoniquement \deleted{\unsure{par des cas}}\inserted{comme un} $(X,S)$
$=(X(I),S(I))$. On a un iso can
